% Experiment 3: Structure A versus Structure B, in the prompt's format, with the
% hidden rules that generate them. Backgrounds are verbatim from held-out prompt
% eval:coin_city:direct_a:r0:k0:cue-correct
% (exp3_training_transfer/coin_city_structural/data/causal/heldout); bold marks the
% phrase that distinguishes them (coloured too, since typewriter fonts often
% lack a bold weight). The table rows are an illustration, not the
% prompt's rows: both cities receive the same news (+4, +12, 0, 0) and support is
% the noise-free worlds.expected_series under that episode's parameters, from
% make_dataset._evaluation_core(COIN_CITY, 0):
%   direct_a   g=0.771169                                   -> 53.1 59.3 50.0 50.0
%   mediated_b g=0.665663 rho=0.550796 phi=0.353853 lambda=0.732322
%                                                           -> 51.9 57.6 56.5 54.4
% Real prompts show 14 weeks with observation noise (sd 1.6).
% The target shares each reference's parameters. The two candidate answers are
% worlds.forecast_scenarios for a +12 shock from rest at horizon 3: A 50.000,
% B 53.647. In worlds.transition the news and
% support in one table row belong to the same week, so the rules are indexed that
% way. The model never sees the rules or parameters. Block titles are relabelled:
% the prompt says REFERENCE SYSTEM 1 / 2 and TARGET SYSTEM, shown here as
% City A / B / C, and the target question is paraphrased from the ten-scenario
% prompt block (the figure shows only its horizon-3 question). Colours follow the transfer-grid figure. Natural width is 7cm
% (half a text width); include it at natural size so the type stays at \scriptsize.
\definecolor{spInk}{HTML}{3F4550}%
\definecolor{spMute}{HTML}{8A93A3}%
\definecolor{spA}{HTML}{4F7EA8}%
\definecolor{spAFill}{HTML}{E8F2FB}%
\definecolor{spB}{HTML}{3F8E82}%
\definecolor{spBFill}{HTML}{E4F2EF}%
\definecolor{spHi}{HTML}{FFF2CC}%
\definecolor{spHiEdge}{HTML}{D6B656}%
\begin{tikzpicture}[
  x=1cm, y=1cm,
  head/.style={anchor=west, font=\ttfamily\bfseries\scriptsize},
  chip/.style={anchor=east, rounded corners=1.6pt, draw=none, text=white,
    inner xsep=3pt, inner ysep=1.4pt, font=\sffamily\bfseries\scriptsize},
  bg/.style={anchor=north west, font=\ttfamily\tiny, text=spInk,
    text width=6.70cm, align=left, inner sep=0pt},
  mono/.style={font=\ttfamily\scriptsize, text=spInk, inner sep=0pt},
  colh/.style={font=\ttfamily\tiny, text=spMute, inner sep=0pt},
  note/.style={anchor=west, font=\sffamily\scriptsize, align=left, inner sep=0pt},
  rule/.style={anchor=north west, font=\scriptsize, text=spInk, inner sep=0pt,
    align=left},
  rulehead/.style={anchor=north west, font=\sffamily\tiny, text=spMute,
    inner sep=0pt},
]
% #1 top y, #2 colour, #3 fill, #4 city letter, #5 chip, #6 background,
% #7 rows (week/news/support), #8 card height
\newcommand{\spcard}[8]{%
  \begin{scope}[shift={(0,#1)}]
    \fill[#3, rounded corners=3pt] (0,0) rectangle (7.0,-#8);
    \draw[#2, rounded corners=3pt, line width=.6pt] (0,0) rectangle (7.0,-#8);
    \node[head, text=#2] at (0.15,-0.25) {CITY #4};
    \node[chip, fill=#2] at (6.85,-0.25) {#5};
    \node[bg] at (0.15,-0.47) {Background: #6};
    \node[colh, anchor=east] at (0.72,-1.30) {week};
    \node[colh, anchor=east] at (2.07,-1.30) {net news};
    \node[colh, anchor=east] at (3.37,-1.30) {support};
    \draw[spMute!60, line width=.3pt] (0.15,-1.41) -- (3.52,-1.41);
    % the two zero-news weeks
    \fill[spHi] (0.09,-1.41-0.28*3+0.125) rectangle (3.55,-1.41-0.28*4-0.135);
    \draw[spHiEdge, line width=.4pt]
      (0.09,-1.41-0.28*3+0.125) rectangle (3.55,-1.41-0.28*4-0.135);
    \foreach \r/\n/\s in {#7}{%
      \node[mono, anchor=east] at (0.72,-1.41-0.28*\r) {\r};
      \node[mono, anchor=east] at (2.07,-1.41-0.28*\r) {\n};
      \node[mono, anchor=east] at (3.37,-1.41-0.28*\r) {\s};}
  \end{scope}}
% hidden-rule panel: #1 top y (absolute), #2 bottom y, #3 colour
\newcommand{\sprule}[3]{%
  \fill[white, rounded corners=2pt] (0.15,#1) rectangle (6.85,#2);
  \draw[#3!45, rounded corners=2pt, line width=.4pt] (0.15,#1) rectangle (6.85,#2);
  \node[rulehead] at (0.28,#1-0.08) {HIDDEN RULE \ (never shown to the model)};}

% ---- City A: Structure A --------------------------------------------------
\spcard{0}{spA}{spAFill}{A}{Structure A: direct}
  {Residents receive campaign developments through synchronized citywide
   alerts. Reactions occur mostly {\color{spA}\bfseries during the same week as the news}.}
  {1/4.0/53.1, 2/12.0/59.3, 3/0.0/50.0, 4/0.0/50.0}
  {3.82}
\node[note, text=spA] at (3.80,-2.43) {no news:\\back at 50};
\draw[spHiEdge, line width=.5pt] (3.55,-2.43) -- (3.72,-2.43);
\sprule{-2.72}{-3.67}{spA}
\node[rule] at (0.28,-3.02) {%
  $y_t = 50 + g\,x_t$\\[1pt]
  {\color{spMute}$g=0.771$;\ \sffamily only this week's news counts}};

% ---- City B: Structure B --------------------------------------------------
\spcard{-3.97}{spB}{spBFill}{B}{Structure B: persistent}
  {Campaign developments usually spread through neighborhood conversations and
   local organizations. Reactions may {\color{spB}\bfseries build gradually and persist} after
   the original story has faded.}
  {1/4.0/51.9, 2/12.0/57.6, 3/0.0/56.5, 4/0.0/54.4}
  {4.14}
\node[note, text=spB] at (3.80,-3.97-2.43) {no news, yet\\still above 54};
\draw[spHiEdge, line width=.5pt] (3.55,-3.97-2.43) -- (3.72,-3.97-2.43);
\sprule{-3.97-2.72}{-3.97-3.99}{spB}
\node[rule] at (0.28,-3.97-3.02) {%
  $m_t = \rho\,m_{t-1} + g\,x_t$
  \quad {\color{spMute}\sffamily($m$: hidden local state)}\\[1pt]
  $y_t - 50 = \phi\,(y_{t-1}-50) + \lambda\,m_t$\\[1pt]
  {\color{spMute}$g=0.666,\ \rho=0.551,\ \phi=0.354,\ \lambda=0.732$}};

% ---- City C: the target, and the two candidate answers ---------------------
\begin{scope}[shift={(0,-8.26)}]
\fill[black!3, rounded corners=3pt] (0,0) rectangle (7.0,-1.95);
\draw[spMute!60, rounded corners=3pt, line width=.6pt] (0,0) rectangle (7.0,-1.95);
\node[head, text=spInk] at (0.15,-0.25) {CITY C};
\node[anchor=west, font=\ttfamily\tiny, text=spInk, inner sep=0pt]
  at (1.28,-0.25) {No observations from City C are available yet.};
\node[bg] at (0.15,-0.47) {Apply net campaign news of +12.0 in the next week,
  then set it to 0. Forecast City C's candidate support 3 weeks later.};
\node[note, text=spA] at (0.15,-1.25)
  {\bfseries If City C works like City A:\ \ predict 50.0};
\node[note, text=spB] at (0.15,-1.65)
  {\bfseries If City C works like City B:\ \ predict 53.6};
\end{scope}
\end{tikzpicture}
